\clearpage
\section{Solutions to the exercises}
\label{sec-solutions-6-finite-fields}

Exact arithmetic and character orthogonality play complementary roles in
these solutions. Root calculations justify modular transforms, while
polynomial and Fourier arguments determine code parameters, Boolean-function
properties, and bounds on distance distributions.

\begin{correction}{exo-calculus-pol-cyclotomic}
For $n>1$, write $r=\operatorname{rad}(n)$ for the product of its distinct
prime factors. Then $\Phi_n(X)=\Phi_r(X^{n/r})$, so repeated prime
factors only insert zero coefficients. For odd $r>1$,
$\Phi_{2r}(X)=\Phi_r(-X)$, which preserves absolute values. Thus a first
occurrence of an absolute value greater than one must have an odd
squarefree index. The value one first occurs at $n=1$, since
$\Phi_1=X-1$; there is no cyclotomic polynomial $\Phi_0$ in this convention.
For distinct primes $p,q$, the coefficients of $\Phi_{pq}$ also have
absolute value at most one. Indeed,
\[
 \Phi_{pq}(X)=\frac{(1-X^{pq})(1-X)}{(1-X^p)(1-X^q)}.
\]
Below degree $pq$, its coefficient at $k$ is $N(k)-N(k-1)$, where
$N(k)$ counts nonnegative solutions of $up+vq=k$. There is at most
one such solution when $k<pq$: two would differ by $(q,-p)$ times
a nonzero integer and force one representation to have value at
least $pq$. Since $\deg\Phi_{pq}<pq$, all coefficients belong to
$\{-1,0,1\}$. Prime indices give coefficients equal to one.

Program~\ref{listing-calcul-pol-cyclo} can therefore skip indices with
fewer than three distinct prime factors when seeking magnitudes above
one. The exact identity
$\Phi_n(X)=\prod_{d\mid n}(1-X^d)^{\mu(n/d)}$ for $n>1$
allows efficient coefficient computation: multiply or divide by
$1-X^d$ in a power series truncated at degree $\varphi(n)$.
The two loop directions in the program implement these operations
without symbolic polynomial expansion.

Searching through $20615$ gives the first indices for $a=1,\ldots,27$:
\[
\begin{array}{c|rrrrrrrrr}
a&1&2&3&4&5&6&7&8&9\\\hline
n&1&105&385&1365&1785&2805&3135&6545&6545\\[2pt]
a&10&11&12&13&14&15&16&17&18\\\hline
n&10465&10465&10465&10465&10465&11305&11305&11305&11305\\[2pt]
a&19&20&21&22&23&24&25&26&27\\\hline
n&11305&11305&11305&15015&11305&17255&17255&20615&20615
\end{array}
\]
These values need not increase: the first index for magnitude $23$
is smaller than that for magnitude $22$. The explicit search limit
also distinguishes an incomplete search from a proof of nonexistence.
\begin{listing}[H]
\PythonListing{finite_fields.py}{14}{30}
\PythonListing{finite_fields.py}{33}{52}

\caption{First occurrences of cyclotomic coefficient magnitudes.}
\label{listing-calcul-pol-cyclo}
\end{listing}
\end{correction}

\begin{correction}{exo-demo-cns-existance-ntt}
\begin{enumerate}
\item Over a field, every nonzero element is regular, so principality
means exact order $n$. The cyclic group $\FF_p^*$ has an element of that
order precisely when $n\mid p-1$.
\item The integer $a$ is prime to $p$, hence invertible modulo $p^r$.
As $a^n=1+pc$ for some integer $c$, repeated use of
$(1+p^ju)^p\equiv1\pmod{p^{j+1}}$ gives
\[
 \zeta_0^n=(1+pc)^{p^{r-1}}\equiv1\pmod{p^r}.
\]
Frobenius is the identity on $\FF_p$, so
$a^{p^{r-1}}\equiv a\pmod p$. These are separate requirements: merely
showing $\zeta_0^{p-1}=1$ would not establish its order divides $n$.
\item For $1\leq s<n$, $a^s\ne1\pmod p$. Thus
$1-\zeta_0^s$ is prime to $p$ and therefore a unit modulo $p^r$.
This proves principality.
\item If $n\mid p_i-1$ for every $i$, combine the lifted roots by the
Chinese remainder theorem. Their differences from lower powers of one
are units in every component, hence units modulo $m$. Conversely, in a
finite ring a regular element is a unit. Reducing a principal root modulo
each $p_i$ therefore preserves all the nonvanishing differences
$1-\zeta^s$, so its order remains exactly $n$. Thus every $p_i-1$ is
divisible by $n$.
\end{enumerate}
\end{correction}

\begin{correction}{exo-root-square-primitive}
\begin{enumerate}
\item If $\zeta$ is principal of order $nm$, the powers in the question
are principal by direct substitution into the definition. Conversely,
the two assumptions imply $\zeta^{nm}=1$. Suppose $a\zeta^i=a$ with
$1\leq i<nm$, and write $i=nq+r$, $0\leq r<n$. If $r=0$, then
$1\leq q<m$, so principality of $\zeta^n$ gives $a=0$. If $r>0$,
raising the equality through $m$ successive multiplications gives
$a(\zeta^m)^r=a$; principality of $\zeta^m$ again gives $a=0$.
\item If $\zeta^2=1$ and $\zeta-1$ is regular, then
$(\zeta-1)(\zeta+1)=0$ forces $\zeta=-1$. This root is principal
exactly when $-2$ is regular. A general ring can have other square roots
of unity, but they are not principal in this sense.
\item Necessity follows by taking the $2^{k-1}$st power and applying the
preceding part. For sufficiency, use induction. When $k>1$, $\zeta^2$
is a principal $2^{k-1}$st root by induction, while
$\zeta^{2^{k-1}}=-1$ is a principal square root. Part~1 concludes.
\end{enumerate}
\end{correction}

\begin{correction}{exo-boolean-functions}
\begin{enumerate}
\item The point indicator at $u\in\FF_2^n$ is the multilinear polynomial
\[
 \delta_u(x)=\prod_{j:u_j=1}x_j\prod_{j:u_j=0}(1+x_j).
\]
Thus such polynomials span all functions. There are $2^n$ monomials,
so both the polynomial space and the function space have dimension
$2^n$ over $\FF_2$; the representation is unique. Each space contains
$2^{2^n}$ elements. Distance zero from the affine subspace means exactly
that the function is affine.
\item The signed correlation with the affine function $(a,b)$ is
\[
 \sum_x(-1)^{\widetilde f(x)+\dotp{a}{x}+b}
    =(-1)^b\Ww f(a)=2^n-2d(\widetilde f,\widetilde f_{a,b}).
\]
Optimize first over $b$, then over all $a$, including zero. The FWT
computes these correlations in $O(n2^n)$ operations.
\item Parseval gives $\sum_a|\Ww f(a)|^2=2^{2n}$. Hence the largest
squared coefficient is at least $2^n$, with equality precisely when
all squared magnitudes equal $2^n$. Walsh coefficients are integers,
so the magnitude $2^{n/2}$ is possible only for even $n$.
\item Since $h(u,v)=f(u)g(v)$, its transform factors as
$\Ww h(a,b)=\Ww f(a)\Ww g(b)$. Flat spectra give a flat product.
Conversely, if the product is flat and nonzero, fixing either index
shows that each factor has constant magnitude; Parseval fixes their
magnitudes to the bent values. The truth table of $u_1u_2$ has signed
Walsh transform $(2,2,2,-2)$. Consequently
$u_1u_2+u_3u_4+\cdots+u_{n-1}u_n$ is bent for even $n$.
\item The truth-table code has length $2^n$, dimension $n+1$, and
$2^{n+1}$ words. A nonconstant affine function is balanced; the
constant-one function has weight $2^n$. Thus its minimum distance is
$2^{n-1}$. If $a$ is stored as a binary integer, the Walsh transform of
$\delta_a$ is the sign table of $\dotp{a}{x}$. Program~\ref{listing-encode-rm}
converts this to bits and adds $b$. For decoding, transform the received
sign table, maximize absolute correlation over every index, and use the
sign to choose $b$. Ties correspond to multiple nearest codewords;
the program chooses the first maximizing index.
\begin{listing}[H]
\PythonListing{error_correcting_codes.py}{15}{21}

\caption{Encoding a first-order Reed--Muller word.}
\label{listing-encode-rm}
\end{listing}
\begin{listing}[H]
\PythonListing{error_correcting_codes.py}{24}{32}

\caption{Nearest-word Reed--Muller decoding by Walsh correlation.}
\label{listing-decode-rm}
\end{listing}
\end{enumerate}
\end{correction}

\begin{correction}{exo-learning-boolean-functions}
\begin{enumerate}
\item Orthogonality of the Walsh basis and $\mathbb E[f^2]=1$ give
$\mathbb E[(f-a_\beta\chi_\beta)^2]=1-a_\beta^2$.
\item If the predicted sign is wrong, $f$ and $h$ lie on opposite
sides of zero, or $h=0$ and $f=-1$. In either case $|f-h|\geq1$.
Thus the error indicator is bounded pointwise by $(f-h)^2$; take expectations.
\item Apply Hoeffding to $Z_j=f(X_j)\chi_\beta(X_j)$, whose mean is
$a_\beta$. The sample-size condition makes the failure probability at
most $\delta$. Conditional on $|\widetilde a_\beta-a_\beta|\leq\lambda$,
orthogonal projection gives
\[
 \mathbb E_X[(f-\widetilde a_\beta\chi_\beta)^2]
   =1-a_\beta^2+(a_\beta-\widetilde a_\beta)^2
   \leq1-a_\beta^2+\lambda^2.
\]
The sign bound then proves the assertion. The fixed-index assumption
matters: choosing an index from the data requires a simultaneous bound.
\item Use $\lambda=\sqrt{\epsilon/M}$. For each of the $M$ retained
coefficients, Hoeffding bounds the probability of an error greater than
$\lambda$ by $\delta/M$. A union bound implies that \emph{all} errors
are at most $\lambda$ with probability at least $1-\delta$; these
coefficient errors need not be independent. On that event, Parseval gives
\[
 \mathbb E[(f-\varphi)^2]=\sum_{w(s)>d}a_s^2+
            \sum_{w(s)\leq d}(a_s-\widetilde a_s)^2
       \leq\alpha+M\lambda^2=\alpha+\epsilon.
\]
Apply part~2. The exact count $M=\sum_{j=0}^d\binom nj$ also covers
$d=0$ and small $n$; replacing it indiscriminately by $n^d$ does not.
\end{enumerate}
\end{correction}

\begin{correction}{exo-matrices-generator-control}
\begin{enumerate}
\item A generator of $\Cc^\bot$ is $H^{\mathsf T}$ and a parity-check
matrix is $G^{\mathsf T}$, since $HG=0$ and both have the required rank.
\item If $G=\begin{pmatrix}I_m\\A\end{pmatrix}$, then
$H=\begin{pmatrix}-A&I_{n-m}\end{pmatrix}$. The first $m$ coordinates
of $Gu$ are the message $u$. The minus sign disappears only in
characteristic two.
\item Choose $m$ independent rows of $G$ and move them to the top by
a coordinate permutation. If the resulting top block is $B$, multiply
$G$ on the right by $B^{-1}$. The permutation gives an equivalent code;
the invertible change of message basis preserves the code itself.
\end{enumerate}
\end{correction}

\begin{correction}{exo-codes-hamming}
\begin{enumerate}
\item The matrix $H$ contains every coordinate vector, so it has rank
$k$ and nullity $2^k-1-k$. No column is zero and no two are equal;
therefore no nonzero codeword has weight one or two. For distinct nonzero
vectors $a,b$, the three distinct columns $a,b,a+b$ sum to zero, giving
minimum distance three. If a single error occurs in position $j$,
its syndrome is column $j$ of $H$. With columns ordered by the binary
integers $1,\ldots,n$, the syndrome directly encodes $j+1$ when positions
are numbered from zero.
\begin{listing}[H]
\PythonListing{error_correcting_codes.py}{35}{38}
\PythonListing{error_correcting_codes.py}{41}{52}

\caption{Syndrome decoding with columns ordered by binary integers.}
\label{listing-decodage-hamming}
\end{listing}
\begin{listing}[H]
\PythonListing{error_correcting_codes.py}{9}{12}

\caption{Binary digits, least significant coordinate first.}
\label{binary-writing-listing}
\end{listing}
\item Let $\alpha$ generate $\FF_{2^k}^*$. A binary polynomial $P$
satisfies $P(\alpha^2)=P(\alpha)^2$, so the BCH conditions for designed
distance three reduce to $P(\alpha)=0$. Expanding
$\sum_{j=0}^{n-1}P_j\alpha^j=0$ in a basis of the extension gives a
binary check matrix whose columns enumerate every nonzero vector of
$\FF_2^k$. This is a Hamming code; requiring all $n$ roots to vanish
would instead give the zero code.
\item Each radius-one ball has $n+1=2^k$ elements. The balls are
disjoint because the minimum distance is three, and their total size
is $2^{n-k}2^k=2^n$. For the coordinate order of
Example~\ref{exmp-code-hamming-7}, a matching matrix is
\[
 H=\begin{pmatrix}1&0&0&1&0&1&1\\0&1&0&1&1&1&0\\0&0&1&0&1&1&1\end{pmatrix}.
\]
Its columns are the coordinates of $1,\alpha,\ldots,\alpha^6$ modulo
$\alpha^3+\alpha+1$, and direct multiplication gives $HG=0$.
It is a column permutation of the canonical binary-integer check matrix.
\item Every nonzero dual word is $(\dotp{u}{v})_{v\ne0}$ for some
$u\ne0$. A nonzero linear functional on $\FF_2^k$ takes each value
$2^{k-1}$ times. Removing the zero input does not remove a one, so its
weight is $2^{k-1}$. Differences of distinct dual words are nonzero
dual words, hence all pairwise distances have this value.
\item Use one nonzero representative from each projective line in
$\FF_q^k$. There are $n=(q^k-1)/(q-1)$ such lines. The check matrix
has rank $k$; no two columns are proportional, while three columns in
a common two-dimensional subspace are dependent. Thus the parameters
are $[n,n-k,3]$. A radius-one ball has $1+n(q-1)=q^k$ elements,
which again proves perfection.
\end{enumerate}
\end{correction}

\begin{correction}{exo-pol-enumerateur-hamming}
\begin{enumerate}
\item Enumeration gives $W_\Cc=X^7+7X^4Y^3+7X^3Y^4+Y^7$.
\item The simplex dual has one zero word and $2^k-1=n$ words of
weight $2^{k-1}=(n+1)/2$. Hence
\[
 W_{\Cc^\bot}=X^n+nX^{(n-1)/2}Y^{(n+1)/2}.
\]
\item Apply the MacWilliams identity to $\Cc^\bot$, whose cardinality
is $n+1$, and use $(\Cc^\bot)^\bot=\Cc$.
\item Expanding at $X=1$ gives
\[
 \begin{aligned}
 A_1=A_2&=0,& A_3&=\frac{n(n-1)}6,\\
 A_4&=\frac{n(n-1)(n-3)}{24},&
 A_5&=\frac{n(n-1)(n-3)(n-7)}{120},\\
 A_6&=\frac{n(n-1)(n-3)(n-5)(n-7)}{720}.
 \end{aligned}
\]
Program~\ref{listing-repartition-hamming} computes the entire enumerator
with exact symbolic arithmetic. Alternatively, each pair of distinct
columns of $H$ has a unique third column equal to their sum. Every
weight-three word is counted by three such pairs, giving
$A_3=\binom n2/3$.
\begin{listing}[H]
\PythonListing{error_correcting_codes.py}{55}{60}

\caption{Exact weight distribution of a binary Hamming code.}
\label{listing-repartition-hamming}
\end{listing}
\item Each row of $H$ contains $2^{k-1}$ ones, an even number for
$k\geq2$. Thus $H\mathbf1=0$, and $x\mapsto x+\mathbf1$ bijects words
of weight $i$ with words of weight $n-i$.
\end{enumerate}
\end{correction}

\begin{correction}{exo-code-hamming-extended}
Appending parity maps weights three and four to four, and weight seven
to eight. The code therefore has parameters $[8,4,4]$ and enumerator
$X^8+14X^4Y^4+Y^8$. Its words, in the order of the messages used in
Example~\ref{exmp-code-hamming-7}, are
\[
 \begin{array}{rrrr}
00000000&11010001&01101001&10111000\\
00110101&11100100&01011100&10001101\\
00011011&11001010&01110010&10100011\\
00101110&11111111&01000111&10010110
\end{array}

\]
Each column of the generator $G$ in that example has weight three, so
the extended generator is $G_8=\begin{pmatrix}G\\1\;1\;1\;1\end{pmatrix}$.
Directly, $G_8^{\mathsf T}G_8=0$ over $\FF_2$. Its rank is four, so its
image is self-orthogonal of dimension half the length and is therefore
self-dual. The enumerator's invariance is a consequence of this argument;
by itself it would only show that the code and its dual have the same
weight distribution.
\end{correction}

\begin{correction}{exo-code-repetition}
\begin{enumerate}
\item The code has dimension one, minimum distance $n$, and enumerator
$W_\Cc=X^n+(q-1)Y^n$.
\item The dual is $\{x:\sum_i x_i=0\}$, of dimension $n-1$ and
minimum distance two for $n\geq2$: $(1,-1,0,\ldots,0)$ is a word,
whereas no weight-one word satisfies the check. Its enumerator is
\[
 W_{\Cc^\bot}=\frac{(X+(q-1)Y)^n+(q-1)(X-Y)^n}{q}.
\]
Self-duality first forces $1=n-1$, hence $n=2$. In that case it holds
exactly when $(1,1)$ is orthogonal to itself, namely in characteristic two.
\end{enumerate}
\end{correction}

\begin{correction}{exo-codes-hadamard}
\begin{enumerate}
\item For the Walsh matrix of order $2^k$, the rows after conversion
are the truth tables of the linear forms $x\mapsto\dotp{a}{x}$.
Removing the coordinate $x=0$ gives the simplex code of dimension $k$.
Adding complements supplies all affine forms and gives the punctured
Reed--Muller code of dimension $k+1$. Both are linear. Linearity is not
guaranteed for a general Hadamard matrix: for order twelve, the two
cardinalities are twelve and twenty-four, neither a power of two.
One Paley construction gives the following twelve words and their
complements:
\[
\begin{array}{c|c}
 \Aa_{12}&\Bb_{12}\setminus\Aa_{12}\\\hline
 00000000000&11111111111\\11011100010&00100011101\\
 01101110001&10010001110\\10110111000&01001000111\\
 01011011100&10100100011\\00101101110&11010010001\\
 00010110111&11101001000\\10001011011&01110100100\\
 11000101101&00111010010\\11100010110&00011101001\\
 01110001011&10001110100\\10111000101&01000111010
\end{array}.
\]
\item Orthogonal sign rows agree in $n/2$ positions and disagree in
$n/2$. The removed first position was common to all rows, so the code
$\Aa_n$ has length $n-1$, cardinality $n$, and all distinct-word distances
$n/2$. For distinct $u,v\in\Aa_n$,
\[
 d(u,\overline v)=(n-1)-d(u,v)=n/2-1,
 \qquad d(u,\overline u)=n-1.
\]
The original and complemented sets are disjoint, since $n/2<n-1$ for
$n\geq4$. Thus $\Bb_n$ has length $n-1$, cardinality $2n$, and minimum
distance $n/2-1$.
\end{enumerate}
\end{correction}

\begin{correction}{exo-polynomials-mds}
\begin{enumerate}
\item The coefficient of $Y^k$ in $(1+Y)^{n-i}(1-Y)^i$ is $P_k(i)$.
Comparing coefficients in $2^mW_{\Cc^\bot}(1,Y)=W_\Cc(1+Y,1-Y)$
gives the formula.
\item Start with
$2^{n-m}W_\Cc(X,1)=\sum_i A'_i(X+1)^{n-i}(X-1)^i$.
After $k$ derivatives, the left side at $X=1$ is
$2^{n-m}k!\sum_{i=0}^{n-k}\binom{n-i}{k}A_i$.
For the summand indexed by $i$ on the right, only the Leibniz term
that differentiates $(X-1)^i$ exactly $i$ times can survive at $X=1$.
For $i\leq k$ its value is
\[
 k!\binom{n-i}{k-i}2^{n-k}A'_i.
\]
Divide by $2^{n-m}k!$ and use
$\binom{n-i}{k-i}=\binom{n-i}{n-k}$.
\item Let $H$ be a full-rank $(n-m)\times n$ check matrix of $\Cc$.
The MDS distance means that every $n-m$ columns are independent.
A nonzero word of the dual has the form $u^{\mathsf T}H$, with $u\ne0$.
If it had $n-m$ zero coordinates, $u$ would annihilate an invertible
submatrix of $H$, a contradiction. Its weight is therefore at least
$m+1$, and Singleton's bound gives equality. In the moment identity
for $k<m$, only $A'_0=1$ survives on the right; on the left, isolate
$A_0=1$ and use $A_1=\cdots=A_{d-1}=0$.
\item Take the equations in the order $k=m-1,m-2,\ldots,0$. The first
contains only $A_d$; each following equation introduces one new
coefficient, with coefficient one. They determine $A_d,\ldots,A_n$
successively. The first gives
$A_d=\binom n{m-1}=\binom nd$.
\end{enumerate}
\end{correction}

\begin{correction}{exo-terminal-linear-programming}
\begin{enumerate}
\item The coefficient comparison in the preceding exercise applies to
$W_{D_\Cc}$ as well: $|\Cc|B'_k=\sum_i B_iP_k(i)$. Proposition~\ref{prop-positivite-repartition-duale}
makes this nonnegative.
\item For minimum distance at least $d$, $B_1,\ldots,B_{d-1}$ vanish,
$B_0=1$, all other $B_i$ are nonnegative, and $\sum_iB_i=|\Cc|$.
Every code therefore supplies a feasible point with objective value its
cardinality, proving the upper bound.
To see boundedness, put $z_k=\sum_i x_iP_k(i)\geq0$. The generating
identity $\sum_kP_k(i)Y^k=(1+Y)^{n-i}(1-Y)^i$ at $Y=1$ gives
\[
 \sum_k z_k=2^n x_0=2^n.
\]
Since $z_0=\sum_i x_i$ and the $z_k$ are nonnegative, the objective
is at most $2^n$. Each nonnegative $x_i$ is bounded by the same value.
The feasible set is closed and nonempty, so the maximum is attained.
\end{enumerate}
\end{correction}
